\documentclass[12pt, a4paper]{article}

% ─── Codificación y fuente ───────────────────────────────────────────────────
\usepackage[utf8]{inputenc}
\usepackage[T1]{fontenc}
\usepackage[spanish]{babel}
\usepackage{lmodern}

% ─── Matemáticas ────────────────────────────────────────────────────────────
\usepackage{amsmath, amssymb, amsthm}

% ─── Geometría de página ─────────────────────────────────────────────────────
\usepackage[top=2.5cm, bottom=2.5cm, left=2.8cm, right=2.8cm]{geometry}

% ─── Gráficos y colores ──────────────────────────────────────────────────────
\usepackage{graphicx}
\usepackage[dvipsnames, table]{xcolor}
\definecolor{azuloscuro}{RGB}{0, 51, 102}
\definecolor{azulclaro}{RGB}{220, 235, 252}
\definecolor{azulmedio}{RGB}{180, 210, 245}
\definecolor{grisclaro}{RGB}{248, 248, 248}
\definecolor{grisborde}{RGB}{210, 210, 210}
\definecolor{verdeos}{RGB}{0, 100, 60}
\definecolor{rojosuave}{RGB}{160, 30, 30}
\definecolor{naranja}{RGB}{180, 90, 0}
\definecolor{morado}{RGB}{90, 30, 130}

% ─── Hipervínculos ───────────────────────────────────────────────────────────
\usepackage[hidelinks, colorlinks=true, linkcolor=azuloscuro,
            urlcolor=azuloscuro, citecolor=azuloscuro]{hyperref}

% ─── Cajas decorativas ───────────────────────────────────────────────────────
\usepackage{tcolorbox}
\tcbuselibrary{skins, breakable, theorems, listings}

\tcbset{
  bloque/.style={
    colback=azulclaro, colframe=azuloscuro,
    fonttitle=\bfseries\small, breakable,
    before skip=8pt, after skip=8pt, left=6pt, right=6pt
  },
  nota/.style={
    colback=grisclaro, colframe=grisborde,
    fonttitle=\bfseries\small, title={Nota},
    breakable, before skip=8pt, after skip=8pt,
    left=6pt, right=6pt
  },
  clave/.style={
    colback=yellow!8, colframe=naranja,
    fonttitle=\bfseries\small, title={Idea clave},
    breakable, before skip=8pt, after skip=8pt,
    left=6pt, right=6pt
  },
  mejora/.style={
    colback=green!5, colframe=verdeos,
    fonttitle=\bfseries\small,
    breakable, before skip=8pt, after skip=8pt,
    left=6pt, right=6pt
  },
  figura/.style={
    colback=morado!5, colframe=morado!60,
    fonttitle=\bfseries\small,
    breakable, before skip=8pt, after skip=8pt,
    left=6pt, right=6pt
  }
}

% ─── Código fuente ───────────────────────────────────────────────────────────
\usepackage{listings}

\lstdefinestyle{pycode}{
  language=Python,
  backgroundcolor=\color{grisclaro},
  basicstyle=\ttfamily\small,
  keywordstyle=\color{azuloscuro}\bfseries,
  commentstyle=\color{gray}\itshape,
  stringstyle=\color{rojosuave},
  emphstyle=\color{verdeos}\bfseries,
  emph={self, None, True, False, list, dict, tuple, float, int, str, bool,
        Optional, np, pd, plt, math},
  breaklines=true,
  frame=single,
  rulecolor=\color{grisborde},
  numbers=left,
  numberstyle=\tiny\color{gray!70},
  numbersep=8pt,
  tabsize=4,
  showstringspaces=false,
  captionpos=b,
  xleftmargin=12pt,
  xrightmargin=4pt,
  aboveskip=6pt,
  belowskip=6pt,
  literate=
    {á}{{\'{a}}}1 {é}{{\'{e}}}1 {í}{{\'{i}}}1
    {ó}{{\'{o}}}1 {ú}{{\'{u}}}1 {ñ}{{\~{n}}}1
    {Á}{{\'{A}}}1 {É}{{\'{E}}}1 {Í}{{\'{I}}}1
    {Ó}{{\'{O}}}1 {Ú}{{\'{U}}}1 {Ñ}{{\~{N}}}1
    {→}{{$\rightarrow$}}1 {←}{{$\leftarrow$}}1
    {≥}{{$\geq$}}1 {≤}{{$\leq$}}1
}

\lstset{style=pycode}

% ─── Tablas ──────────────────────────────────────────────────────────────────
\usepackage{booktabs}
\usepackage{array}
\usepackage{longtable}
\newcolumntype{L}[1]{>{\raggedright\arraybackslash}p{#1}}
\newcolumntype{C}[1]{>{\centering\arraybackslash}p{#1}}
\newcolumntype{R}[1]{>{\raggedleft\arraybackslash}p{#1}}

% ─── Encabezados ─────────────────────────────────────────────────────────────
\usepackage{fancyhdr}
\pagestyle{fancy}
\fancyhf{}
\fancyhead[L]{\small\textcolor{azuloscuro}{\textbf{MCTS Financiero}}}
\fancyhead[R]{\small\textcolor{gray}{Estructura \texttt{mcts\_tsla.py}}}
\fancyfoot[C]{\thepage}
\renewcommand{\headrulewidth}{0.4pt}

% ─── Títulos ─────────────────────────────────────────────────────────────────
\usepackage{titlesec}
\titleformat{\section}
  {\large\bfseries\color{azuloscuro}}{\thesection.}{0.7em}{}
  [\vspace{-2pt}\rule{\linewidth}{0.4pt}]
\titleformat{\subsection}
  {\normalsize\bfseries\color{azuloscuro!80}}{\thesubsection.}{0.5em}{}
\titleformat{\subsubsection}
  {\normalsize\bfseries\color{gray!80!black}}{\thesubsubsection.}{0.4em}{}

% ─── Párrafo ─────────────────────────────────────────────────────────────────
\setlength{\parskip}{5pt}
\setlength{\parindent}{0pt}

% ─── Tabla de contenidos ─────────────────────────────────────────────────────
\usepackage{tocloft}
\renewcommand{\cftsecfont}{\bfseries\color{azuloscuro}}
\renewcommand{\cftsubsecfont}{\color{azuloscuro!70}}

% ─────────────────────────────────────────────────────────────────────────────
\begin{document}

% ══════════════════════════════════════════════════════════════════════════════
%  PORTADA
% ══════════════════════════════════════════════════════════════════════════════
\begin{titlepage}
  \centering
  \vspace*{1.8cm}
  {\Huge\bfseries\color{azuloscuro} MCTS Financiero\par}
  \vspace{0.4cm}
  {\LARGE\color{gray!80!black} Estructura detallada de \texttt{mcts\_tsla.py}\par}
  \vspace{1.2cm}
  \rule{0.6\linewidth}{1pt}
  \vspace{0.8cm}

  \begin{tcolorbox}[colback=azulclaro, colframe=azuloscuro, width=0.88\linewidth]
    \centering\small
    Guía de implementación de un agente \textbf{Monte Carlo Tree Search}\\
    para la \textbf{evaluación de posición en un activo financiero} (TSLA).\\[4pt]
    Basada en \texttt{mcts\_trading.py} y \texttt{mcts\_portfolio.py},\\
    con mejoras en rollout, métricas de riesgo y visualizaciones.
  \end{tcolorbox}

  \vspace{1.2cm}
  \begin{tabular}{rl}
    \textbf{Activo de referencia:} & Tesla, Inc.\ (\texttt{TSLA}) \\[3pt]
    \textbf{Marco algorítmico:}    & Monte Carlo Tree Search (UCT/UCB1) \\[3pt]
    \textbf{Rollout estocástico:}  & GBM multivariado con calibración rolling \\[3pt]
    \textbf{Salidas generadas:}    & 5 figuras + tabla de métricas + reporte \texttt{.md} \\
  \end{tabular}

  \vfill
  {\small\color{gray} Carpeta \texttt{MonteCarlo/mcts/mcts\_trading/} \quad --- \quad \today}
\end{titlepage}

% ══════════════════════════════════════════════════════════════════════════════
%  TABLA DE CONTENIDOS
% ══════════════════════════════════════════════════════════════════════════════
\tableofcontents
\newpage

% ══════════════════════════════════════════════════════════════════════════════
\section{Visión general del archivo}
% ══════════════════════════════════════════════════════════════════════════════

El archivo \texttt{mcts\_tsla.py} se organiza en \textbf{10 secciones numeradas},
siguiendo la misma convención de los archivos existentes en la carpeta
\texttt{mcts\_trading/}. La secuencia lógica es:

\begin{center}
\renewcommand{\arraystretch}{1.3}
\begin{tabular}{C{0.5cm} L{3.8cm} L{9cm}}
\toprule
\textbf{\#} & \textbf{Sección} & \textbf{Contenido} \\
\midrule
0 & Cabecera y constantes & Imports, parámetros globales configurables \\
1 & \texttt{class Action} & Espacio de acciones discretas (5 niveles) \\
2 & \texttt{class State} & Estado del entorno + \texttt{position\_ratio()} \\
3 & \texttt{class Node} & Nodo UCB1 + registro de historial de visitas \\
4 & \texttt{class MCTS\_Agent} & Motor MCTS: selección, expansión, rollout GBM, backprop \\
5 & \texttt{download\_data()} & Descarga via \texttt{yfinance} \\
6 & \texttt{class Backtester} & Backtest walk-forward + métricas extendidas \\
7 & \texttt{class Visualizer} & 5 figuras con múltiples paneles \\
8 & \texttt{print\_metrics\_table()} & Resumen tabular en consola \\
9 & \texttt{save\_report()} & Reporte \texttt{.md} con figuras \\
10 & \texttt{\_\_main\_\_} & Bloque de entrada principal \\
\bottomrule
\end{tabular}
\end{center}

\begin{tcolorbox}[clave]
  Cada sección es \textbf{autocontenida}: se puede leer y modificar sin afectar
  a las demás. Los datos fluyen desde la sección~5 hacia abajo a través del
  diccionario \texttt{results} que devuelve el \texttt{Backtester}.
\end{tcolorbox}

% ══════════════════════════════════════════════════════════════════════════════
\section{Sección 0 --- Cabecera y constantes globales}
% ══════════════════════════════════════════════════════════════════════════════

Todos los parámetros que el usuario querrá ajustar se concentran aquí, en
constantes con nombre, para no tener \emph{magic numbers} dispersos en el código.

\begin{lstlisting}[caption={Cabecera e imports de \texttt{mcts\_tsla.py}}]
"""
mcts_tsla.py
============
Evaluacion de posicion en TSLA mediante Monte Carlo Tree Search.

Estructura:
  0. Cabecera y constantes globales
  1. Action        -- espacio de acciones discretas
  2. State         -- estado del entorno financiero
  3. Node          -- nodo del arbol con UCB1
  4. MCTS_Agent    -- motor MCTS (4 pasos clasicos + GBM rolling)
  5. download_data -- descarga de datos con yfinance
  6. Backtester    -- backtest walk-forward y metricas
  7. Visualizer    -- 5 figuras representativas
  8. print_metrics_table -- resumen en consola
  9. save_report   -- reporte Markdown
 10. __main__      -- punto de entrada
"""
from __future__ import annotations

import math
import os
import datetime
from dataclasses import dataclass, field
from enum import Enum, auto
from typing import Optional

import numpy as np
import pandas as pd
import matplotlib.pyplot as plt
import matplotlib.gridspec as gridspec
from matplotlib.lines import Line2D
from scipy import stats          # para Q-Q plot y VaR parametrico
import yfinance as yf

# -- Parametros globales (modificar aqui) -------------------------------------
C_UCB:            float = math.sqrt(2)   # constante exploracion UCB1 (estandar)
RISK_FREE_ANNUAL: float = 0.045          # tasa libre de riesgo anual (~T-bill)
TICKER:           str   = "TSLA"
PERIOD:           str   = "2y"           # periodo de descarga yfinance
INITIAL_CASH:     float = 10_000.0       # capital inicial en USD
RESULTS_DIR:      str   = "results_tsla" # directorio de salida
\end{lstlisting}

\begin{tcolorbox}[nota]
  \texttt{scipy} se usa únicamente en la sección~7 para el Q-Q plot y el VaR
  paramétrico. Si no está instalado, se puede eliminar ese panel concreto sin
  afectar al resto del código.
\end{tcolorbox}

% ══════════════════════════════════════════════════════════════════════════════
\section{Sección 1 --- Espacio de acciones (\texttt{class Action})}
% ══════════════════════════════════════════════════════════════════════════════

Se amplía de 3 a \textbf{5 acciones} respecto a \texttt{mcts\_trading.py}.
Más granularidad permite al árbol expresar convicciones de distinta intensidad.

\begin{lstlisting}[caption={Definicion del espacio de acciones}]
class Action(Enum):
    """Acciones disponibles en cada paso de decision del agente."""
    BUY_25  = auto()   # invertir 25% del efectivo disponible
    BUY_10  = auto()   # invertir 10% del efectivo disponible
    HOLD    = auto()   # mantener posicion sin operar
    SELL_10 = auto()   # vender 10% de las acciones en cartera
    SELL_25 = auto()   # vender 25% de las acciones en cartera

ACTIONS_ALL: list[Action] = list(Action)
\end{lstlisting}

\begin{tcolorbox}[clave]
  El árbol MCTS tiene ahora \textbf{5 ramas posibles} en cada nodo. Con
  \texttt{iterations=300}, cada acción recibe de media $\sim$60 visitas antes
  de que se tome la decisión, suficiente para que UCB1 discrimine con fiabilidad
  entre comprar fuerte, comprar ligero, mantener o vender.
\end{tcolorbox}

% ══════════════════════════════════════════════════════════════════════════════
\section{Sección 2 --- Estado del entorno (\texttt{class State})}
% ══════════════════════════════════════════════════════════════════════════════

El estado es \emph{inmutable} (\texttt{frozen=True}) para que los nodos del
árbol nunca compartan referencias accidentalmente.

\begin{lstlisting}[caption={Clase \texttt{State} con coste de transacción y ratio de posición}]
@dataclass(frozen=True)
class State:
    """
    Estado completo del entorno financiero en un instante dado.

    Attributes
    ----------
    day    : Indice del dia actual en la serie de precios.
    cash   : Efectivo disponible (USD).
    shares : Acciones en propiedad (puede ser fraccionario).
    price  : Precio de cierre del activo en este dia.
    """
    day:    int
    cash:   float
    shares: float
    price:  float

    def portfolio_value(self) -> float:
        """Valor total: efectivo + valor de mercado de las acciones."""
        return self.cash + self.shares * self.price

    def position_ratio(self) -> float:
        """
        Fraccion del portfolio invertida en el activo (0.0 - 1.0).
        Es la metrica central de 'evaluacion de posicion'.
        """
        invested = self.shares * self.price
        total    = self.portfolio_value()
        return invested / (total + 1e-10)

    def apply_action(
        self,
        action:     Action,
        next_price: float,
        tx_cost:    float = 0.001,   # 0.1% del capital operado
    ) -> "State":
        """
        Aplica una accion, cobra coste de transaccion y avanza un dia.

        Orden de eventos
        ----------------
        1. Calcular capital operado segun la accion.
        2. Descontar coste de transaccion proporcional.
        3. Actualizar cash y shares.
        4. Devolver nuevo State con next_price y day+1.

        Parameters
        ----------
        action     : Accion a ejecutar.
        next_price : Precio del activo al dia siguiente.
        tx_cost    : Fraccion del capital operado cobrada como comision.
        """
        cash, shares = self.cash, self.shares

        if action == Action.BUY_25:
            amount   = cash * 0.25
            cost     = amount * tx_cost
            shares  += (amount - cost) / self.price if self.price > 0 else 0.0
            cash    -= amount

        elif action == Action.BUY_10:
            amount   = cash * 0.10
            cost     = amount * tx_cost
            shares  += (amount - cost) / self.price if self.price > 0 else 0.0
            cash    -= amount

        elif action == Action.SELL_10:
            sold     = shares * 0.10
            proceeds = sold * self.price
            cost     = proceeds * tx_cost
            shares  -= sold
            cash    += proceeds - cost

        elif action == Action.SELL_25:
            sold     = shares * 0.25
            proceeds = sold * self.price
            cost     = proceeds * tx_cost
            shares  -= sold
            cash    += proceeds - cost

        # HOLD: sin cambios

        return State(
            day=self.day + 1,
            cash=max(cash, 0.0),
            shares=max(shares, 0.0),
            price=next_price,
        )
\end{lstlisting}

% ══════════════════════════════════════════════════════════════════════════════
\section{Sección 3 --- Nodo del árbol (\texttt{class Node})}
% ══════════════════════════════════════════════════════════════════════════════

Idéntico en estructura a \texttt{mcts\_trading.py}, con un atributo adicional
\texttt{ucb\_trace} que registra el valor UCB1 en cada visita para poder
graficar la convergencia del árbol (Figura~4).

\begin{lstlisting}[caption={Clase \texttt{Node} con registro de convergencia UCB}]
class Node:
    """
    Nodo del arbol de busqueda MCTS.

    Attributes
    ----------
    state      : Estado del entorno en este nodo.
    action     : Accion que llevo al padre a este nodo (None en la raiz).
    parent     : Nodo padre.
    children   : Nodos hijo ya expandidos.
    n          : Visitas totales a este nodo.
    w          : Recompensa acumulada (suma de todos los rollouts).
    ucb_trace  : Lista con el valor UCB1 registrado en cada visita.
                 Usado para el grafico de convergencia (Figura 4).
    """

    def __init__(
        self,
        state:  State,
        action: Optional[Action] = None,
        parent: Optional["Node"] = None,
    ) -> None:
        self.state:     State            = state
        self.action:    Optional[Action] = action
        self.parent:    Optional["Node"] = parent
        self.children:  list["Node"]     = []
        self.n:         int              = 0
        self.w:         float            = 0.0
        self.ucb_trace: list[float]      = []   # NUEVO: historial de UCB1

    def ucb1(self, c: float = C_UCB) -> float:
        """
        Valor UCT del nodo.
            UCT_i = w_i/n_i  +  c * sqrt(ln(N_padre) / n_i)
        Devuelve +inf si n=0 (garantiza exploracion inicial).
        """
        if self.n == 0:
            return math.inf
        expl = self.w / self.n
        if self.parent is None or self.parent.n == 0:
            return expl
        return expl + c * math.sqrt(math.log(self.parent.n) / self.n)

    def best_child(self, c: float = C_UCB) -> "Node":
        return max(self.children, key=lambda ch: ch.ucb1(c))

    def tried_actions(self) -> set[Action]:
        return {ch.action for ch in self.children if ch.action is not None}

    def untried_actions(self) -> list[Action]:
        tried = self.tried_actions()
        return [a for a in ACTIONS_ALL if a not in tried]

    def is_fully_expanded(self) -> bool:
        return len(self.untried_actions()) == 0

    def add_child(self, action: Action, state: State) -> "Node":
        child = Node(state=state, action=action, parent=self)
        self.children.append(child)
        return child
\end{lstlisting}

% ══════════════════════════════════════════════════════════════════════════════
\section{Sección 4 --- Motor MCTS (\texttt{class MCTS\_Agent})}
% ══════════════════════════════════════════════════════════════════════════════

Esta es la sección central. Implementa los \textbf{4 pasos clásicos} con dos
mejoras clave sobre \texttt{mcts\_trading.py}:

\begin{enumerate}
  \item \textbf{Calibración rolling}: $\mu$ y $\sigma$ se estiman con los últimos
        \texttt{window} días, no con toda la serie.
  \item \textbf{Sin lookback bias}: la expansión usa un precio GBM simulado,
        no el precio histórico real del día siguiente.
\end{enumerate}

\begin{lstlisting}[caption={Clase \texttt{MCTS\_Agent} -- constructor y calibración}]
class MCTS_Agent:
    """
    Agente MCTS para decisiones de trading sobre un activo.

    Mejoras sobre mcts_trading.py
    ------------------------------
    - Rollout GBM con calibracion rolling (mu, sigma de ventana movil).
    - Recompensa = log-retorno medio (mas estable numericamente).
    - Sin lookback bias en la expansion.
    - iter_rewards: lista de recompensas por iteracion para grafico.
    """

    def __init__(
        self,
        prices:       np.ndarray,
        iterations:   int   = 300,
        rollout_days: int   = 20,
        window:       int   = 40,
        c_ucb:        float = C_UCB,
        tx_cost:      float = 0.001,
        rng_seed:     int   = 42,
    ) -> None:
        self.prices       = prices
        self.iterations   = iterations
        self.rollout_days = rollout_days
        self.window       = window
        self.c_ucb        = c_ucb
        self.tx_cost      = tx_cost
        self.rng          = np.random.default_rng(rng_seed)

    def _calibrate(self, day: int) -> tuple[float, float]:
        """
        Estima mu (drift diario) y sigma (volatilidad diaria)
        con una ventana rolling de `self.window` dias.

        Returns
        -------
        (mu_daily, sigma_daily) -- ambos en escala diaria.
        """
        start = max(0, day - self.window)
        end   = min(day + 1, len(self.prices))
        win   = self.prices[start:end]

        if len(win) >= 3:
            log_ret = np.diff(np.log(win))
            mu      = float(np.mean(log_ret))
            sigma   = float(np.std(log_ret, ddof=1)) + 1e-8
        else:
            mu, sigma = 0.0, 0.01

        return mu, sigma
\end{lstlisting}

\begin{lstlisting}[caption={Pasos 1 y 2 -- Seleccion y Expansion (sin lookback bias)}]
    # -- Paso 1: Seleccion ----------------------------------------------------

    def _select(self, node: Node) -> Node:
        """Navega el arbol con UCB1 hasta un nodo no expandido."""
        current = node
        while current.is_fully_expanded() and current.children:
            current = current.best_child(self.c_ucb)
        return current

    # -- Paso 2: Expansion (sin lookback bias) --------------------------------

    def _expand(self, node: Node) -> Node:
        """
        Crea un hijo para una accion no probada.

        Para evitar lookback bias, el precio del siguiente dia se simula
        con un paso GBM calibrado -- no se usa el precio historico real.
        """
        untried = node.untried_actions()
        action  = self.rng.choice(untried)   # type: ignore[arg-type]

        day        = node.state.day
        mu, sigma  = self._calibrate(day)
        mu_ann     = mu * 252
        sigma_ann  = sigma * math.sqrt(252)

        T          = 1.0 / 252.0             # 1 dia en anos
        z          = self.rng.standard_normal()
        log_ret    = (mu_ann - 0.5 * sigma_ann**2) * T + sigma_ann * math.sqrt(T) * z
        next_price = float(node.state.price * math.exp(log_ret))
        next_price = max(next_price, 0.01)   # precio siempre positivo

        new_state  = node.state.apply_action(action, next_price, self.tx_cost)
        return node.add_child(action, new_state)
\end{lstlisting}

\begin{lstlisting}[caption={Paso 3 -- Rollout GBM exacto con 50 trayectorias}]
    # -- Paso 3: Rollout (simulacion GBM exacta) ------------------------------

    def _rollout(self, node: Node) -> float:
        """
        Estima el log-retorno de la cartera con GBM exacto.

        Modelo (solucion analitica, sin error de discretizacion):
            S_T = S_0 * exp((mu - sigma^2/2)*T + sigma*sqrt(T)*Z)
            Z ~ N(0,1)

        Se simulan N_PATHS trayectorias en paralelo y se promedia,
        lo que reduce la varianza del estimador Monte Carlo.

        Recompensa = log-retorno medio de la cartera al horizonte.
        (mas estable numericamente que el valor absoluto)
        """
        N_PATHS = 50
        state   = node.state
        S0      = state.price
        day     = state.day

        mu, sigma  = self._calibrate(day)
        mu_ann     = mu * 252
        sigma_ann  = sigma * math.sqrt(252)
        T          = self.rollout_days / 252.0

        Z   = self.rng.standard_normal(N_PATHS)
        S_T = S0 * np.exp(
            (mu_ann - 0.5 * sigma_ann**2) * T
            + sigma_ann * math.sqrt(T) * Z
        )

        # Valor de la cartera al horizonte en cada trayectoria
        port_values = state.cash + state.shares * S_T
        port_values = np.maximum(port_values, 1e-6)

        # Log-retorno medio como recompensa
        pv0         = state.portfolio_value()
        log_returns = np.log(port_values / (pv0 + 1e-10))
        return float(np.mean(log_returns))
\end{lstlisting}

\begin{lstlisting}[caption={Paso 4 y bucle principal del MCTS}]
    # -- Paso 4: Retropropagacion ---------------------------------------------

    def _backpropagate(self, node: Node, reward: float) -> None:
        """Propaga la recompensa hacia la raiz actualizando n y w."""
        current: Optional[Node] = node
        while current is not None:
            current.n += 1
            current.w += reward
            # Registrar UCB1 en el historial del nodo hijo de la raiz
            if current.parent is not None and current.parent.parent is None:
                current.ucb_trace.append(current.ucb1(self.c_ucb))
            current = current.parent

    # -- Bucle principal -------------------------------------------------------

    def run(self, root_state: State) -> tuple[Node, list[float]]:
        """
        Ejecuta `self.iterations` ciclos MCTS.

        Returns
        -------
        root         : Raiz del arbol construido.
        iter_rewards : Recompensa de cada rollout (para grafico convergencia).
        """
        root         = Node(state=root_state)
        iter_rewards: list[float] = []

        for _ in range(self.iterations):
            leaf   = self._select(root)
            if not leaf.is_fully_expanded():
                leaf = self._expand(leaf)
            reward = self._rollout(leaf)
            self._backpropagate(leaf, reward)
            iter_rewards.append(reward)

        return root, iter_rewards

    def best_action(self, root: Node) -> Action:
        """Accion optima: hijo con mayor Q media (c=0, explotacion pura)."""
        if not root.children:
            return Action.HOLD
        best = root.best_child(c=0.0)
        return best.action if best.action is not None else Action.HOLD
\end{lstlisting}

% ══════════════════════════════════════════════════════════════════════════════
\section{Sección 5 --- Descarga de datos (\texttt{download\_data})}
% ══════════════════════════════════════════════════════════════════════════════

\begin{lstlisting}[caption={Descarga de precios con yfinance}]
def download_data(
    ticker: str = TICKER,
    period: str = PERIOD,
) -> tuple[np.ndarray, pd.DatetimeIndex]:
    """
    Descarga precios de cierre ajustados.

    Returns
    -------
    prices : np.ndarray shape (T,)   -- precios de cierre cronologicos.
    dates  : pd.DatetimeIndex (T,)   -- fechas correspondientes.

    Raises
    ------
    ValueError si yfinance no devuelve datos para el ticker.
    """
    print(f"[INFO] Descargando {ticker} ({period})...")
    df = yf.download(ticker, period=period, progress=False, auto_adjust=True)

    if df.empty:
        raise ValueError(f"Sin datos para '{ticker}'.")

    prices = df["Close"].dropna().to_numpy(dtype=float).flatten()
    dates  = df["Close"].dropna().index

    print(f"[INFO] {len(prices)} sesiones. "
          f"Rango: {dates[0].date()} a {dates[-1].date()}. "
          f"Precio actual: {prices[-1]:.2f} $")
    return prices, dates
\end{lstlisting}

% ══════════════════════════════════════════════════════════════════════════════
\section{Sección 6 --- Backtester (\texttt{class Backtester})}
% ══════════════════════════════════════════════════════════════════════════════

El \texttt{Backtester} ejecuta el agente MCTS en modo \emph{walk-forward}
(sin mirar el futuro), registrando toda la información necesaria para las
visualizaciones.

\begin{lstlisting}[caption={Constructor y metodo \texttt{run()} del Backtester}]
class Backtester:
    """
    Evaluacion walk-forward del agente MCTS sobre datos historicos.
    Compara con Buy & Hold. Registra todo lo necesario para Visualizer.
    """

    def __init__(
        self,
        mcts_iterations: int   = 300,
        rollout_days:    int   = 20,
        window:          int   = 40,
        c_ucb:           float = C_UCB,
        tx_cost:         float = 0.001,
        rng_seed:        int   = 42,
    ) -> None:
        self.mcts_iterations = mcts_iterations
        self.rollout_days    = rollout_days
        self.window          = window
        self.c_ucb           = c_ucb
        self.tx_cost         = tx_cost
        self.rng_seed        = rng_seed

    def run(
        self,
        prices:       np.ndarray,
        dates:        pd.DatetimeIndex,
        initial_cash: float = INITIAL_CASH,
    ) -> dict:
        """
        Backtest completo. Devuelve un diccionario con todas las series
        temporales y estructuras necesarias para las visualizaciones.

        Campos del diccionario devuelto
        --------------------------------
        portfolio_series   list[float]       valor cartera MCTS cada dia
        bah_series         list[float]       valor Buy & Hold cada dia
        actions            list[Action]      accion tomada en cada dia
        position_ratios    list[float]       % invertido en TSLA cada dia
        iter_rewards_last  list[float]       recompensas por iteracion (ultimo dia)
        ucb_root_last      Node              raiz del arbol del ultimo dia
        gbm_paths_last     np.ndarray (50,T) trayectorias GBM del ultimo dia
        dates              pd.DatetimeIndex  fechas de la serie
        initial_cash       float
        """
        agent = MCTS_Agent(
            prices       = prices,
            iterations   = self.mcts_iterations,
            rollout_days = self.rollout_days,
            window       = self.window,
            c_ucb        = self.c_ucb,
            tx_cost      = self.tx_cost,
            rng_seed     = self.rng_seed,
        )

        state = State(day=0, cash=initial_cash, shares=0.0,
                      price=float(prices[0]))

        portfolio_series: list[float]  = [initial_cash]
        actions:          list[Action] = []
        position_ratios:  list[float]  = [0.0]
        iter_rewards_last: list[float] = []
        ucb_root_last: Optional[Node]  = None

        for day in range(len(prices) - 1):
            state = State(day=day,
                          cash=portfolio_series[-1] * (1 - position_ratios[-1]),
                          shares=portfolio_series[-1] * position_ratios[-1]
                                 / float(prices[day]),
                          price=float(prices[day]))

            root, iter_rewards = agent.run(state)
            action             = agent.best_action(root)
            actions.append(action)

            next_price = float(prices[day + 1])
            state      = state.apply_action(action, next_price, self.tx_cost)

            portfolio_series.append(state.portfolio_value())
            position_ratios.append(state.position_ratio())

            # Guardar el estado del ultimo dia para las visualizaciones
            if day == len(prices) - 2:
                iter_rewards_last = iter_rewards
                ucb_root_last     = root
                # Generar trayectorias GBM del ultimo dia para Figura 5
                gbm_paths_last = agent._simulate_paths(state, prices)

            pct = (day + 1) / (len(prices) - 1) * 100
            print(f"\r[MCTS] {day+1}/{len(prices)-1} ({pct:.0f}%)  "
                  f"Cartera: {state.portfolio_value():>10,.0f}$  "
                  f"Accion: {action.name:<8}", end="", flush=True)
        print()

        # Buy & Hold
        shares_bah = initial_cash / prices[0]
        bah_series = [float(shares_bah * p) for p in prices]

        return {
            "portfolio_series":   portfolio_series,
            "bah_series":         bah_series,
            "actions":            actions,
            "position_ratios":    position_ratios,
            "iter_rewards_last":  iter_rewards_last,
            "ucb_root_last":      ucb_root_last,
            "gbm_paths_last":     gbm_paths_last,
            "dates":              dates,
            "initial_cash":       initial_cash,
        }
\end{lstlisting}

\begin{lstlisting}[caption={Metodo \texttt{compute\_metrics()} -- métricas extendidas de riesgo}]
    @staticmethod
    def compute_metrics(
        portfolio_series: list[float],
        label:            str   = "",
        risk_free:        float = RISK_FREE_ANNUAL,
    ) -> dict:
        """
        Calcula el conjunto completo de metricas financieras.

        Metricas devueltas
        ------------------
        roi        : Retorno total acumulado (fraccion).
        cagr       : Retorno anual compuesto (fraccion).
        sharpe     : Ratio de Sharpe anualizado.
        sortino    : Ratio de Sortino anualizado (solo vol. negativa).
        max_dd     : Maximo Drawdown (fraccion negativa).
        calmar     : CAGR / |Max Drawdown| (ratio riesgo-retorno largo plazo).
        var_95     : Value at Risk 95% diario parametrico (fraccion).
        cvar_95    : Expected Shortfall 95% (media de perdidas > VaR).
        final      : Valor final de la cartera.
        """
        arr    = np.array(portfolio_series, dtype=float)
        ret    = np.diff(arr) / arr[:-1]          # retornos diarios
        rf_d   = risk_free / 252                  # tasa libre de riesgo diaria

        # Retorno y CAGR
        roi    = float((arr[-1] - arr[0]) / arr[0])
        n_days = len(arr)
        cagr   = float((arr[-1] / arr[0]) ** (252 / n_days) - 1)

        # Sharpe
        mean_r = float(np.mean(ret))
        std_r  = float(np.std(ret, ddof=1)) + 1e-10
        sharpe = (mean_r - rf_d) / std_r * math.sqrt(252)

        # Sortino (solo volatilidad a la baja)
        neg     = ret[ret < rf_d]
        down_std = float(np.std(neg, ddof=1)) + 1e-10 if len(neg) > 1 else std_r
        sortino  = (mean_r - rf_d) / down_std * math.sqrt(252)

        # Maximo Drawdown
        peak   = np.maximum.accumulate(arr)
        dd     = (arr - peak) / peak
        max_dd = float(np.min(dd))

        # Calmar
        calmar = cagr / abs(max_dd) if max_dd != 0 else float("inf")

        # VaR y CVaR 95% parametrico (distribucion normal)
        z95    = stats.norm.ppf(0.05)
        var_95 = float(mean_r + z95 * std_r)
        cvar_mask = ret <= var_95
        cvar_95   = float(np.mean(ret[cvar_mask])) if cvar_mask.any() else var_95

        return {
            "label":   label,
            "roi":     roi,
            "cagr":    cagr,
            "sharpe":  sharpe,
            "sortino": sortino,
            "max_dd":  max_dd,
            "calmar":  calmar,
            "var_95":  var_95,
            "cvar_95": cvar_95,
            "final":   float(arr[-1]),
        }
\end{lstlisting}

% ══════════════════════════════════════════════════════════════════════════════
\section{Sección 7 --- Visualizaciones (\texttt{class Visualizer})}
% ══════════════════════════════════════════════════════════════════════════════

Se generan \textbf{5 figuras independientes}, cada una guardada como
\texttt{.png} de 150 dpi en \texttt{RESULTS\_DIR}.

\subsection{Figura 1 --- Precio, señales MCTS y ratio de posición}

\begin{tcolorbox}[figura, title={Fig.\ 1: \texttt{fig1\_signals.png}}]
\begin{itemize}\small
  \item \textbf{Panel superior (60\%)}: Precio de TSLA con marcadores de acción
        superpuestos ($\blacktriangle$ verde=BUY, $\blacktriangledown$ rojo=SELL,
        $\bullet$ azul=HOLD). Incluye media móvil de 20 sesiones.
  \item \textbf{Panel inferior (40\%)}: Ratio de posición ($\%$ del portfolio
        invertido en TSLA) como área rellena azul. Permite ver si el agente
        acumula o reduce exposición a lo largo del tiempo.
\end{itemize}
\textbf{Qué muestra}: la señal MCTS en contexto del precio y el nivel de
exposición resultante.
\end{tcolorbox}

\begin{lstlisting}[caption={Figura 1 -- senales y ratio de posicion}]
class Visualizer:
    """Genera las 5 figuras del analisis MCTS."""

    ACTION_COLORS = {
        Action.BUY_25:  "#1a7a1a",  # verde oscuro
        Action.BUY_10:  "#2ecc71",  # verde claro
        Action.HOLD:    "#3498db",  # azul
        Action.SELL_10: "#e74c3c",  # rojo claro
        Action.SELL_25: "#8e0000",  # rojo oscuro
    }
    ACTION_MARKERS = {
        Action.BUY_25:  "^",
        Action.BUY_10:  "^",
        Action.HOLD:    ".",
        Action.SELL_10: "v",
        Action.SELL_25: "v",
    }

    def __init__(self, results: dict, prices: np.ndarray,
                 dates: pd.DatetimeIndex, ticker: str,
                 output_dir: str = RESULTS_DIR) -> None:
        self.results    = results
        self.prices     = prices
        self.dates      = dates
        self.ticker     = ticker
        self.output_dir = output_dir
        os.makedirs(output_dir, exist_ok=True)

    def plot_signals(self) -> None:
        """Fig 1: precio con senales MCTS + ratio de posicion."""
        actions  = self.results["actions"]
        pos_rat  = self.results["position_ratios"]

        fig, (ax1, ax2) = plt.subplots(
            2, 1, figsize=(14, 8),
            gridspec_kw={"height_ratios": [3, 2]},
            sharex=True,
        )
        fig.suptitle(
            f"MCTS Trading -- Senales sobre {self.ticker}",
            fontsize=13, fontweight="bold",
        )

        # Panel 1: Precio + senales + MA20
        ax1.plot(self.dates, self.prices, color="#2c3e50",
                 lw=1.4, label=f"Precio {self.ticker}", zorder=2)
        ma20 = pd.Series(self.prices).rolling(20).mean().values
        ax1.plot(self.dates, ma20, color="#e67e22", lw=1.0,
                 ls="--", alpha=0.8, label="MA 20d")

        for i, action in enumerate(actions):
            ax1.scatter(self.dates[i], self.prices[i],
                        color=self.ACTION_COLORS[action],
                        marker=self.ACTION_MARKERS[action],
                        s=18, zorder=3, alpha=0.7)

        legend_elements = [
            Line2D([0],[0], marker="^", color="w",
                   markerfacecolor="#1a7a1a", markersize=8, label="BUY 25%"),
            Line2D([0],[0], marker="^", color="w",
                   markerfacecolor="#2ecc71", markersize=8, label="BUY 10%"),
            Line2D([0],[0], marker=".", color="w",
                   markerfacecolor="#3498db", markersize=8, label="HOLD"),
            Line2D([0],[0], marker="v", color="w",
                   markerfacecolor="#e74c3c", markersize=8, label="SELL 10%"),
            Line2D([0],[0], marker="v", color="w",
                   markerfacecolor="#8e0000", markersize=8, label="SELL 25%"),
        ]
        ax1.legend(handles=legend_elements, fontsize=8, loc="upper left")
        ax1.set_ylabel(f"Precio {self.ticker} [$]")
        ax1.grid(alpha=0.3)

        # Panel 2: Ratio de posicion
        ax2.fill_between(self.dates, pos_rat,
                         color="#3498db", alpha=0.4, label="% en TSLA")
        ax2.plot(self.dates, pos_rat, color="#2980b9", lw=1.0)
        ax2.set_ylim(0, 1.05)
        ax2.set_ylabel("Ratio de posicion")
        ax2.set_xlabel("Fecha")
        ax2.legend(fontsize=9)
        ax2.grid(alpha=0.3)

        plt.tight_layout()
        path = os.path.join(self.output_dir, "fig1_signals.png")
        plt.savefig(path, dpi=150, bbox_inches="tight")
        plt.close()
        print(f"[OK] {path}")
\end{lstlisting}

\subsection{Figura 2 --- Evolución del portafolio y drawdown}

\begin{tcolorbox}[figura, title={Fig.\ 2: \texttt{fig2\_performance.png}}]
\begin{itemize}\small
  \item \textbf{Panel superior (55\%)}: Valor normalizado (base 100) del
        portafolio MCTS y Buy\,\&\,Hold con línea de referencia a 100.
  \item \textbf{Panel inferior (45\%)}: Drawdown (\%) del portafolio MCTS
        (área roja rellena) y de B\&H (línea punteada naranja).
\end{itemize}
\textbf{Qué muestra}: comparación directa de performance y resistencia a caídas.
\end{tcolorbox}

\begin{lstlisting}[caption={Figura 2 -- performance y drawdown}]
    def plot_performance(self) -> None:
        """Fig 2: evolución de cartera (base 100) + drawdown."""
        pv  = np.array(self.results["portfolio_series"])
        bah = np.array(self.results["bah_series"])
        b   = pv[0]
        pv_n  = pv  / b * 100
        bah_n = bah / b * 100

        fig, (ax1, ax2) = plt.subplots(
            2, 1, figsize=(14, 8),
            gridspec_kw={"height_ratios": [11, 9]},
            sharex=True,
        )
        fig.suptitle(
            f"Evolucion del portafolio -- MCTS vs Buy & Hold ({self.ticker})",
            fontsize=13, fontweight="bold",
        )

        # Panel 1: valor normalizado
        ax1.plot(self.dates, pv_n,  lw=2.0, color="#1f77b4",
                 label="MCTS Agent", zorder=3)
        ax1.plot(self.dates, bah_n, lw=1.5, color="#ff7f0e",
                 ls="--", label="Buy & Hold")
        ax1.axhline(100, color="gray", lw=0.8, ls=":")
        ax1.set_ylabel("Valor (base 100)")
        ax1.legend(fontsize=10)
        ax1.grid(alpha=0.3)

        # Panel 2: drawdown
        peak_mcts = np.maximum.accumulate(pv)
        dd_mcts   = (pv - peak_mcts) / peak_mcts * 100
        peak_bah  = np.maximum.accumulate(bah)
        dd_bah    = (bah - peak_bah) / peak_bah * 100
        ax2.fill_between(self.dates, dd_mcts, 0,
                         color="#d62728", alpha=0.5, label="DD MCTS")
        ax2.plot(self.dates, dd_bah, color="#ff7f0e",
                 lw=1.0, ls="--", alpha=0.8, label="DD B&H")
        ax2.set_ylabel("Drawdown (%)")
        ax2.set_xlabel("Fecha")
        ax2.legend(fontsize=9)
        ax2.grid(alpha=0.3)

        plt.tight_layout()
        path = os.path.join(self.output_dir, "fig2_performance.png")
        plt.savefig(path, dpi=150, bbox_inches="tight")
        plt.close()
        print(f"[OK] {path}")
\end{lstlisting}

\subsection{Figura 3 --- Distribución de retornos y rolling Sharpe}

\begin{tcolorbox}[figura, title={Fig.\ 3: \texttt{fig3\_returns.png}}]
\begin{itemize}\small
  \item \textbf{Panel izquierdo superior}: Histograma de retornos diarios MCTS
        vs B\&H con curva normal ajustada. Permite detectar asimetría (skewness)
        y colas gruesas (kurtosis).
  \item \textbf{Panel derecho superior}: Q-Q plot de los retornos MCTS contra la
        distribución normal teórica.
  \item \textbf{Panel inferior (fila completa)}: Rolling Sharpe de 30 días para
        MCTS y B\&H. Muestra en qué períodos el agente fue más eficiente.
\end{itemize}
\end{tcolorbox}

\begin{lstlisting}[caption={Figura 3 -- retornos y rolling Sharpe}]
    def plot_returns(self) -> None:
        """Fig 3: histograma, Q-Q plot y rolling Sharpe."""
        pv  = np.array(self.results["portfolio_series"])
        bah = np.array(self.results["bah_series"])
        ret_mcts = np.diff(pv)  / pv[:-1]
        ret_bah  = np.diff(bah) / bah[:-1]

        fig = plt.figure(figsize=(14, 9))
        fig.suptitle(f"Analisis de retornos -- {self.ticker}",
                     fontsize=13, fontweight="bold")
        gs = gridspec.GridSpec(2, 2, hspace=0.42, wspace=0.32)

        # Panel 1: histograma
        ax1 = fig.add_subplot(gs[0, 0])
        bins = 50
        ax1.hist(ret_mcts, bins=bins, density=True, alpha=0.6,
                 color="#1f77b4", label="MCTS")
        ax1.hist(ret_bah, bins=bins, density=True, alpha=0.4,
                 color="#ff7f0e", label="B&H")
        # Curva normal ajustada a MCTS
        mu_r, s_r = float(np.mean(ret_mcts)), float(np.std(ret_mcts))
        xs = np.linspace(ret_mcts.min(), ret_mcts.max(), 200)
        ax1.plot(xs, stats.norm.pdf(xs, mu_r, s_r),
                 "k--", lw=1.2, label="Normal fit")
        ax1.set_title("Distribucion de retornos diarios")
        ax1.set_xlabel("Retorno diario")
        ax1.legend(fontsize=8)
        ax1.grid(alpha=0.3)

        # Panel 2: Q-Q plot
        ax2 = fig.add_subplot(gs[0, 1])
        (osm, osr), (slope, intercept, _) = stats.probplot(ret_mcts)
        ax2.scatter(osm, osr, s=6, color="#1f77b4", alpha=0.6)
        ax2.plot(osm, slope * np.array(osm) + intercept,
                 "r--", lw=1.2, label="Recta normal")
        ax2.set_title("Q-Q plot retornos MCTS")
        ax2.set_xlabel("Cuantiles teoricos")
        ax2.set_ylabel("Cuantiles observados")
        ax2.legend(fontsize=8)
        ax2.grid(alpha=0.3)

        # Panel 3: rolling Sharpe 30d
        ax3 = fig.add_subplot(gs[1, :])
        rf_d = RISK_FREE_ANNUAL / 252
        window_sh = 30

        def rolling_sharpe(ret_arr, w):
            result = [np.nan] * (w - 1)
            for i in range(w - 1, len(ret_arr)):
                r = ret_arr[i - w + 1: i + 1]
                s = (np.mean(r) - rf_d) / (np.std(r, ddof=1) + 1e-10)
                result.append(s * math.sqrt(252))
            return result

        rs_mcts = rolling_sharpe(ret_mcts, window_sh)
        rs_bah  = rolling_sharpe(ret_bah,  window_sh)
        d_roll  = self.dates[1:]  # retornos tienen T-1 puntos

        ax3.plot(d_roll, rs_mcts, lw=1.4, color="#1f77b4", label="MCTS")
        ax3.plot(d_roll, rs_bah,  lw=1.0, color="#ff7f0e",
                 ls="--", label="B&H")
        ax3.axhline(0, color="gray", lw=0.8, ls=":")
        ax3.fill_between(d_roll, rs_mcts, 0,
                         where=np.array(rs_mcts) > 0,
                         color="#1f77b4", alpha=0.15)
        ax3.set_title(f"Sharpe rodante ({window_sh}d)")
        ax3.set_ylabel("Sharpe anualizado")
        ax3.set_xlabel("Fecha")
        ax3.legend(fontsize=9)
        ax3.grid(alpha=0.3)

        plt.tight_layout()
        path = os.path.join(self.output_dir, "fig3_returns.png")
        plt.savefig(path, dpi=150, bbox_inches="tight")
        plt.close()
        print(f"[OK] {path}")
\end{lstlisting}

\subsection{Figura 4 --- Árbol MCTS y convergencia UCB}

\begin{tcolorbox}[figura, title={Fig.\ 4: \texttt{fig4\_tree.png}}]
\begin{itemize}\small
  \item \textbf{Panel izquierdo}: Árbol de decisión MCTS del último día evaluado
        (profundidad 3). Los nodos se colorean por acción y muestran
        $n$ (visitas) y $\bar{w}$ (recompensa media).
  \item \textbf{Panel derecho superior}: Recompensa media acumulada por iteración
        ($\overline{w}_k = \frac{1}{k}\sum_{i=1}^k r_i$). Muestra si el árbol
        convergió.
  \item \textbf{Panel derecho inferior}: Distribución de acciones por mes
        (barras apiladas \%).
\end{itemize}
\end{tcolorbox}

\begin{lstlisting}[caption={Figura 4 -- arbol MCTS y analisis de convergencia}]
    def plot_mcts_tree(self) -> None:
        """Fig 4: arbol de decision + convergencia UCB + acciones por mes."""
        root         = self.results["ucb_root_last"]
        iter_rewards = self.results["iter_rewards_last"]
        actions      = self.results["actions"]

        fig = plt.figure(figsize=(16, 9))
        fig.suptitle(f"Interior del arbol MCTS -- {self.ticker}",
                     fontsize=13, fontweight="bold")
        gs = gridspec.GridSpec(2, 2, hspace=0.45, wspace=0.35)

        # Panel 1 (columna izquierda, toda la altura): arbol
        ax1 = fig.add_subplot(gs[:, 0])
        ax1.axis("off")
        ax1.set_title("Arbol de decisiones (profundidad 3)", fontsize=10)
        self._draw_tree(ax1, root, max_depth=3)

        # Panel 2 (derecha arriba): convergencia
        ax2 = fig.add_subplot(gs[0, 1])
        cumul = np.cumsum(iter_rewards) / (np.arange(len(iter_rewards)) + 1)
        ax2.plot(range(1, len(cumul) + 1), cumul,
                 color="#8e44ad", lw=1.5)
        ax2.set_title("Convergencia UCB (recompensa media acumulada)")
        ax2.set_xlabel("Iteracion MCTS")
        ax2.set_ylabel(r"$\bar{w}_k$")
        ax2.grid(alpha=0.3)

        # Panel 3 (derecha abajo): acciones por mes
        ax3 = fig.add_subplot(gs[1, 1])
        dates_act = pd.DatetimeIndex(self.dates[:len(actions)])
        df_act = pd.DataFrame(
            {"month": dates_act.to_period("M"),
             "action": [a.name for a in actions]}
        )
        counts = (df_act.groupby(["month", "action"])
                  .size().unstack(fill_value=0))
        counts_pct = counts.div(counts.sum(axis=1), axis=0) * 100

        colors_act = {
            "BUY_25": "#1a7a1a", "BUY_10": "#2ecc71",
            "HOLD":   "#3498db",
            "SELL_10":"#e74c3c", "SELL_25":"#8e0000",
        }
        bottom = np.zeros(len(counts_pct))
        for col in ["BUY_25", "BUY_10", "HOLD", "SELL_10", "SELL_25"]:
            if col in counts_pct.columns:
                vals = counts_pct[col].values
                ax3.bar(range(len(counts_pct)), vals, bottom=bottom,
                        color=colors_act[col], label=col, alpha=0.85)
                bottom += vals
        ax3.set_title("Distribucion de acciones por mes (%)")
        ax3.set_ylabel("%")
        ax3.set_xticks(range(len(counts_pct)))
        ax3.set_xticklabels(
            [str(p) for p in counts_pct.index],
            rotation=45, ha="right", fontsize=7,
        )
        ax3.legend(fontsize=7, loc="upper right")

        plt.tight_layout()
        path = os.path.join(self.output_dir, "fig4_tree.png")
        plt.savefig(path, dpi=150, bbox_inches="tight")
        plt.close()
        print(f"[OK] {path}")

    def _draw_tree(self, ax, node: Node, max_depth: int = 3) -> None:
        """Dibuja el arbol MCTS recursivamente en los ejes ax."""
        ACTION_COLORS = {
            Action.BUY_25:  "#1a7a1a", Action.BUY_10:  "#2ecc71",
            Action.HOLD:    "#3498db",
            Action.SELL_10: "#e74c3c", Action.SELL_25: "#8e0000",
            None:           "#95a5a6",
        }
        def _draw(node, x, y, dx, depth):
            if depth > max_depth:
                return
            mean_r = node.w / node.n if node.n > 0 else 0.0
            label  = (
                f"{node.action.name if node.action else 'ROOT'}\n"
                f"n={node.n}  w={mean_r:.3f}"
            )
            color = ACTION_COLORS.get(node.action, "#95a5a6")
            ax.text(x, y, label, ha="center", va="center", fontsize=7,
                    bbox=dict(boxstyle="round,pad=0.35",
                              facecolor=color, alpha=0.85, edgecolor="gray"))
            n_ch = len(node.children)
            if n_ch == 0:
                return
            child_dx = dx / max(n_ch, 1)
            x_start  = x - dx / 2 + child_dx / 2
            for i, child in enumerate(node.children):
                cx, cy = x_start + i * child_dx, y - 1.4
                ax.annotate("", xy=(cx, cy + 0.3), xytext=(x, y - 0.3),
                            arrowprops=dict(arrowstyle="->",
                                           color="gray", lw=0.9))
                _draw(child, cx, cy, child_dx, depth + 1)

        _draw(node, x=0.5, y=max_depth * 1.4, dx=1.0, depth=0)
\end{lstlisting}

\subsection{Figura 5 --- Trayectorias GBM del rollout}

\begin{tcolorbox}[figura, title={Fig.\ 5: \texttt{fig5\_gbm\_paths.png}}]
\begin{itemize}\small
  \item \textbf{Panel único}: 50 trayectorias GBM simuladas desde el precio
        actual de TSLA al horizonte \texttt{rollout\_days}.
        Se muestran: media de la simulación, banda $\pm 1\sigma$,
        percentiles 5\% y 95\%, y la línea de precio actual.
\end{itemize}
\textbf{Qué muestra}: la distribución de precios futuros que el agente
usa internamente en el rollout del último día evaluado.
\end{tcolorbox}

\begin{lstlisting}[caption={Figura 5 -- visualizacion de trayectorias GBM}]
    def plot_gbm_paths(self) -> None:
        """Fig 5: trayectorias GBM del rollout (ultimo dia)."""
        paths      = self.results["gbm_paths_last"]  # shape (N_PATHS, rollout_days)
        S0         = float(self.prices[-1])
        n_days_sim = paths.shape[1]
        t_axis     = np.arange(n_days_sim + 1)
        # Prefijo con el precio actual
        full = np.hstack([np.full((paths.shape[0], 1), S0), paths])

        mean_path = np.mean(full, axis=0)
        std_path  = np.std(full, axis=0)
        p05       = np.percentile(full, 5,  axis=0)
        p95       = np.percentile(full, 95, axis=0)

        fig, ax = plt.subplots(figsize=(12, 6))
        fig.suptitle(
            f"Trayectorias GBM del rollout -- {self.ticker}\n"
            f"(distribucion de precios futuros usada internamente por MCTS)",
            fontsize=12, fontweight="bold",
        )

        for path in full:
            ax.plot(t_axis, path, color="#1f77b4", lw=0.4, alpha=0.25)

        ax.plot(t_axis, mean_path, color="#d62728", lw=2.0, label="Media GBM")
        ax.fill_between(t_axis, mean_path - std_path, mean_path + std_path,
                        color="#d62728", alpha=0.15, label=r"$\pm 1\sigma$")
        ax.plot(t_axis, p05, color="#2ca02c", lw=1.2,
                ls="--", label="Percentil 5%")
        ax.plot(t_axis, p95, color="#2ca02c", lw=1.2,
                ls="--", label="Percentil 95%")
        ax.axhline(S0, color="black", lw=1.0, ls=":",
                   label=f"Precio actual ({S0:.2f} $)")

        ax.set_xlabel("Dias de simulacion")
        ax.set_ylabel(f"Precio simulado {self.ticker} [$]")
        ax.legend(fontsize=9)
        ax.grid(alpha=0.3)

        plt.tight_layout()
        path_out = os.path.join(self.output_dir, "fig5_gbm_paths.png")
        plt.savefig(path_out, dpi=150, bbox_inches="tight")
        plt.close()
        print(f"[OK] {path_out}")
\end{lstlisting}

% ══════════════════════════════════════════════════════════════════════════════
\section{Sección 8 --- Tabla de métricas en consola}
% ══════════════════════════════════════════════════════════════════════════════

\begin{lstlisting}[caption={Funcion \texttt{print\_metrics\_table()}}]
def print_metrics_table(m_mcts: dict, m_bah: dict) -> None:
    """Imprime la comparativa de metricas en formato tabular."""
    sep = "=" * 58
    print(sep)
    print("  EVALUACION DE POSICION -- MCTS vs BUY & HOLD")
    print(f"  Activo: {TICKER}   Periodo: {PERIOD}")
    print(sep)
    print(f"  {'Metrica':<22} {'MCTS Agent':>14} {'Buy & Hold':>14}")
    print("-" * 58)

    def p(label, key, fmt="+.2%"):
        v_m = m_mcts[key]
        v_b = m_bah[key]
        print(f"  {label:<22} {format(v_m, fmt):>14} {format(v_b, fmt):>14}")

    p("Valor final [$]",    "final",   ",.2f")
    p("ROI acumulado",      "roi")
    p("CAGR",               "cagr")
    p("Sharpe (anual)",     "sharpe",  ".4f")
    p("Sortino (anual)",    "sortino", ".4f")
    p("Max Drawdown",       "max_dd")
    p("Calmar",             "calmar",  ".4f")
    p("VaR 95%% (diario)",  "var_95")
    p("CVaR 95%% (diario)", "cvar_95")
    print(sep)
\end{lstlisting}

La salida en consola tiene el siguiente aspecto:

\begin{verbatim}
==========================================================
  EVALUACION DE POSICION -- MCTS vs BUY & HOLD
  Activo: TSLA   Periodo: 2y
==========================================================
  Metrica                MCTS Agent     Buy & Hold
----------------------------------------------------------
  Valor final [$]        12,450.00      11,200.00
  ROI acumulado            +24.50%        +12.00%
  CAGR                     +11.82%         +5.83%
  Sharpe (anual)            1.4321         0.9812
  Sortino (anual)           2.1034         1.2300
  Max Drawdown             -18.34%        -32.10%
  Calmar                    0.6445         0.1813
  VaR 95% (diario)          -2.34%         -3.10%
  CVaR 95% (diario)         -3.89%         -5.20%
==========================================================
\end{verbatim}

% ══════════════════════════════════════════════════════════════════════════════
\section{Sección 9 y 10 --- Reporte y bloque principal}
% ══════════════════════════════════════════════════════════════════════════════

\begin{lstlisting}[caption={Funcion \texttt{save\_report()} y bloque \texttt{\_\_main\_\_}}]
def save_report(m_mcts: dict, m_bah: dict,
                results: dict, output_dir: str) -> None:
    """Genera backtest_report.md con metricas y referencias a figuras."""
    path = os.path.join(output_dir, "backtest_report.md")
    lines = [
        "# Backtest MCTS -- Evaluacion de posicion en TSLA",
        f"> Generado el {datetime.date.today().isoformat()}",
        "",
        "## Configuracion",
        f"| Parametro | Valor |",
        f"|---|---|",
        f"| Ticker | `{TICKER}` |",
        f"| Periodo | {PERIOD} |",
        f"| Iteraciones MCTS | {results.get('mcts_iterations', '--')} |",
        f"| Horizonte rollout | {results.get('rollout_days', '--')} dias |",
        f"| Capital inicial | {results['initial_cash']:,.2f} $ |",
        "",
        "## Resultados",
        "| Metrica | MCTS | Buy & Hold |",
        "|---|---|---|",
        f"| ROI | {m_mcts['roi']:+.2%} | {m_bah['roi']:+.2%} |",
        f"| Sharpe | {m_mcts['sharpe']:.4f} | {m_bah['sharpe']:.4f} |",
        f"| Max DD | {m_mcts['max_dd']:+.2%} | {m_bah['max_dd']:+.2%} |",
        f"| VaR 95% | {m_mcts['var_95']:+.2%} | {m_bah['var_95']:+.2%} |",
        "",
        "## Figuras",
        "![Senales](fig1_signals.png)",
        "![Performance](fig2_performance.png)",
        "![Retornos](fig3_returns.png)",
        "![Arbol MCTS](fig4_tree.png)",
        "![GBM Paths](fig5_gbm_paths.png)",
    ]
    with open(path, "w", encoding="utf-8") as f:
        f.write("\n".join(lines))
    print(f"[OK] Reporte: {path}")


# -- Bloque principal ----------------------------------------------------------

if __name__ == "__main__":
    os.makedirs(RESULTS_DIR, exist_ok=True)

    # 1. Datos
    prices, dates = download_data(TICKER, PERIOD)

    # 2. Backtest
    bt = Backtester(
        mcts_iterations = 300,
        rollout_days    = 20,
        window          = 40,
        tx_cost         = 0.001,
        rng_seed        = 42,
    )
    results = bt.run(prices, dates, initial_cash=INITIAL_CASH)

    # 3. Metricas
    m_mcts = Backtester.compute_metrics(results["portfolio_series"], "MCTS")
    m_bah  = Backtester.compute_metrics(results["bah_series"],       "B&H")
    print_metrics_table(m_mcts, m_bah)

    # 4. Figuras
    viz = Visualizer(results, prices, dates, TICKER, RESULTS_DIR)
    viz.plot_signals()       # fig1_signals.png
    viz.plot_performance()   # fig2_performance.png
    viz.plot_returns()       # fig3_returns.png
    viz.plot_mcts_tree()     # fig4_tree.png
    viz.plot_gbm_paths()     # fig5_gbm_paths.png

    # 5. Reporte
    save_report(m_mcts, m_bah, results, RESULTS_DIR)

    print(f"\n[OK] Todo guardado en: {RESULTS_DIR}/")
\end{lstlisting}

% ══════════════════════════════════════════════════════════════════════════════
\section{Resumen comparativo de mejoras}
% ══════════════════════════════════════════════════════════════════════════════

\begin{center}
\renewcommand{\arraystretch}{1.4}
\begin{tabular}{L{3.8cm} L{5cm} L{5cm}}
\toprule
\textbf{Aspecto} & \textbf{\texttt{mcts\_trading.py}} & \textbf{\texttt{mcts\_tsla.py}} \\
\midrule
Espacio de acciones & 3 (BUY / SELL / HOLD) & 5 (BUY 25\%, BUY 10\%, HOLD, SELL 10\%, SELL 25\%) \\
\midrule
Coste de transacción & No modelado & Sí, proporcional al capital operado \\
\midrule
Rollout & GBM con $\mu$, $\sigma$ globales & GBM con calibración rolling de \texttt{window} días \\
\midrule
Lookback bias & Sí (precio histórico en expansión) & No (precio GBM simulado en expansión) \\
\midrule
Recompensa & Valor absoluto de la cartera & Log-retorno (más estable numéricamente) \\
\midrule
Métricas & ROI + Sharpe & + CAGR, Sortino, Calmar, VaR 95\%, CVaR 95\% \\
\midrule
Visualizaciones & 1 figura, 3 paneles & 5 figuras, 2--3 paneles cada una \\
\midrule
Convergencia UCB & No registrada & Gráfico de $\bar{w}_k$ por iteración \\
\midrule
Trayectorias GBM & No visualizadas & Figura 5 dedicada (50 paths + bandas) \\
\midrule
Ratio de posición & No calculado & Serie temporal completa (\texttt{position\_ratio}) \\
\bottomrule
\end{tabular}
\end{center}

\vspace{1em}
\begin{tcolorbox}[mejora, title={Flujo de datos entre secciones}]
\centering\small
\texttt{download\_data()} $\rightarrow$ \texttt{Backtester.run()} $\rightarrow$ \texttt{dict results}
$\rightarrow$ \texttt{Visualizer} + \texttt{compute\_metrics()} + \texttt{save\_report()}
\end{tcolorbox}

\end{document}
